\subsection{The price process}
\label{sec:expl:price}

Before any model is trained, the data say a great deal about what can and cannot be
predicted.  This section collects the descriptive results that the methodology and the experiments rely on:
what the stream of book events looks like, how often a price moves over each horizon, how
the price level bounds the dynamics, how regular the activity of a market is, what the kind
of contract says, whether the contracts of the same event lead one another, and whether the
wallets that trade carry information.  Every number is computed on the training and
validation days of \cref{sec:data:protocol}, never on the locked test period.

\subsubsection{The event stream}
\label{sec:expl:stream}

An audit of one training day (12 March,
325 million events on 23{,}512 markets) shows how uneven this unit is.  Only 4.2\% of
the events change the mid (1.9\% at extreme prices, 5.2\% at central ones); 55\% add
size and 44\% remove it.  Ninety percent of the events are away from the best quote
(median distance 0.03 on the bid side and 0.10 on the ask side), and a full book has a
median of about 650 bid and 750 ask levels, so most events never reach the ten stored
levels and enter the vector only through the order-flow accumulators.  The event clock
is irregular: 30\% of consecutive events share a timestamp, the median gap is 31~ms and
the 90th percentile 4.3~s.  One 20-event snapshot therefore spans 12~s in median but
196~s at the 90th percentile, and 73\% of the snapshots contain no mid change; a
10-snapshot window spans 2.8 minutes in median and 34 minutes at the 90th percentile,
with 44\% of the windows containing no mid change.  The consequences for the models are
discussed in \cref{sec:exp:nostate,sec:analysis:nostate}.

\subsubsection{How often a price moves}
\label{sec:expl:classes}

The three-class label of \cref{sec:data:labels} is dominated by \textsc{const} at every
horizon below a few hours (\cref{tab:expl:classes,fig:expl:classes}).  Over one minute nine samples in ten do
not move; over fifteen minutes seven in ten; over one hour a little more than half; over
twelve hours one in three.  The two directions are close to balanced at every horizon, with
\textsc{down} slightly more frequent than \textsc{up} at the short ones.  Two consequences
follow.  The first is that the task changes with the horizon: at one minute it is almost
entirely ``will the mid move at all'', and the direction of a move, once it moves, is nearly
a coin toss (\cref{sec:expl:direction}); at twelve hours the direction is the task.  The
second is that any classifier's balanced accuracy at a short horizon depends on where it
places its decision thresholds far more than on what it knows: a model that under-calls the
moves by a few per cent loses several points of balanced accuracy, which is why every
comparison in \cref{ch:horizons} is made after the same post-hoc calibration.

\begin{table}[t]
\centering
\caption{Share of the three classes among the labelled validation samples of 13 March, by
horizon (samples of the stateful models: every market of every block at the block time;
1 h, 6 h and 12 h from the causal-baseline evaluation, 1 and 15 min from the short-horizon
runs).}
\label{tab:expl:classes}
\begin{fittable}\begin{tabular}{lccc}
\toprule
Horizon & \textsc{down} & \textsc{const} & \textsc{up} \\
\midrule
1 min  & 0.057 & 0.891 & 0.052 \\
15 min & 0.164 & 0.690 & 0.146 \\
1 h    & 0.236 & 0.549 & 0.215 \\
6 h    & 0.335 & 0.354 & 0.312 \\
12 h   & 0.377 & 0.275 & 0.348 \\
\bottomrule
\end{tabular}\end{fittable}
\end{table}

\begin{figure}[t]
\centering
\begin{tikzpicture}
\begin{axis}[thesis, ybar stacked, bar width=22pt, ymin=0, ymax=1.02, ylabel={share of the labelled samples},
  symbolic x coords={1 min, 15 min, 1 h, 6 h, 12 h}, xtick=data, enlarge x limits=0.14,
  legend columns=3, legend style={at={(0.5,1.03)}, anchor=south},
  nodes near coords, every node near coord/.append style={font=\scriptsize, color=white}, point meta=explicit symbolic]
\addplot[fill=vizA, draw=white, line width=0.6pt] coordinates {(1 min,0.057)[] (15 min,0.164)[0.16] (1 h,0.236)[0.24] (6 h,0.335)[0.34] (12 h,0.377)[0.38]};
\addplot[fill=vizB, draw=white, line width=0.6pt] coordinates {(1 min,0.891)[0.89] (15 min,0.690)[0.69] (1 h,0.549)[0.55] (6 h,0.354)[0.35] (12 h,0.275)[0.28]};
\addplot[fill=vizC, draw=white, line width=0.6pt] coordinates {(1 min,0.052)[] (15 min,0.146)[0.15] (1 h,0.215)[0.22] (6 h,0.312)[0.31] (12 h,0.348)[0.35]};
\legend{\textsc{down}, \textsc{const}, \textsc{up}}
\end{axis}
\end{tikzpicture}

\caption{Class shares of the validation labels by horizon.  Below one hour the task is
dominated by \textsc{const}; at twelve hours the three classes are nearly balanced.}
\label{fig:expl:classes}
\end{figure}

\subsubsection{The price level}
\label{sec:expl:level}

About 40\% of the samples have a mid below $0.05$ or above $0.95$.  A Polymarket mid price
is a bounded, discrete process observed through a thin book; most of its changes are
one-tick moves produced by an order arriving at or leaving the best quote, they are not
information and they revert, and a market that has not moved in the last hour is usually
illiquid or settled near a boundary and will not move in the next hour either.  At extreme
prices the process is almost purely a one-tick oscillation: the share of events that change
the mid falls from 5.2\% at central prices to 1.9\% there (\cref{sec:expl:stream}).  Every
result of the thesis is therefore reported separately on the \emph{extreme} samples (mid
outside $[0.05, 0.95]$) and the \emph{central} ones, and every learned model is weakest,
relative to the simple rules, at the extremes.  The literature adds that the extremes are
also where prices are systematically biased: contracts bought below ten cents lose money on
average and contracts above ninety cents earn a little \cite{cardozo2026favorite}, the
favourite--longshot bias of \cref{sec:pm:what}, so the price level is not only a bound on
the dynamics but a feature in its own right.

\subsubsection{The activity clock}
\label{sec:expl:activity}

The most regular thing about a market is its own rhythm.  The time since the mid last
changed, the number of mid changes in the last hour and the number of book events in the
last minute are the three quantities that best separate the samples whose mid will move
from those whose mid will not.  Over a one-hour horizon the time since the last change
separates \textsc{const} from a move with an area under the ROC curve of $0.85$; over one
minute the number of changes in the last hour reaches $0.94$, the time since the last change
$0.91$ and the event count of the last minute $0.89$ (\cref{sec:stateless:control}).  No
quantity of the book's shape comes close.  This is the \emph{activity clock}, and it is the
first fact any model of this data has to know: a market that has moved often lately will
move again, one that has been still will stay still.  It is also a fact that the event-clock
representation of \cref{sec:data:tensors} hides, since a window of twenty book events can be
ten seconds or ten minutes old and nothing in the window says which; \cref{ch:stateless,ch:stateful}
measure what this costs.

\subsubsection{The direction of a move}
\label{sec:expl:direction}

Given that the mid moves, which way is far less predictable.  On the samples that moved
within one hour, the sign of ``mid minus its one-hour mean'' gives the direction 57.3\% of
the time, the level-1 imbalance of the book 55.7\%, the microprice 55.5\%, the on-chain
taker flow 52.6\% (on the fifth of the samples with trades), the ten-level imbalance 50.8\%
and the all-level order-flow accumulators of the tensor 50.0--50.5\%.  The direction is in
the first level of the book and in the position of the mid relative to its recent mean,
which is the reversal of \cref{sec:expl:level}, and not in the depth of the book; this is the
finding of \textcite{tran2022beyondbest} and \textcite{kolm2023dofi} on equities reproduced
on a prediction market.  Over one minute the best univariate hit rate is 54.1\% (the level-1
imbalance), every feature is within four points of a coin toss, and the twelve-hour return
is the strongest single feature at twelve hours with a 57.9\% directional hit as a reversal
sign.  The rule that follows from these numbers, ``predict that the price will move back
against its recent change, and no move if it did not move'', is Mean Reversion-$H$, the
reference of every comparison in the thesis (\cref{sec:models:baselines}).

\subsubsection{The kind of contract}
\label{sec:expl:templates}

Questions are templated: 3{,}528 templates among the active markets (``Will the
temperature in $X$ be above $y$ on $d$'', ``Will $X$ win the $Y$ game'', ``Bitcoin Up or
Down, $t$ to $t + 15$ min'').  The template explains 55\% of the between-market variance of
the probability of a move within one hour, from 4\% for House-seat markets to 83\% for
daily-temperature markets, and a ridge regression on a sentence embedding of the question
predicts that probability out of sample with $R^2 = 0.22$, no better than the template
means.  The contract type therefore conditions the dynamics, as \textcite{le2026calibration}
also find for the calibration of political markets, but it does so through a handful of
categories rather than through the words of the question; the usable form of the text is a
category or template indicator, and that is how it enters the models
(\cref{sec:prep:features}).

\subsubsection{Siblings co-move, nobody leads}
\label{sec:expl:siblings}

Contracts of the same event, the two sides of a game or the candidates of an election, are
the obvious place to look for one price leading another (\cref{fig:expl:leadlag}).  Taking the minute changes of the mid of every market
(6.65 million changes over 49 million market-minutes of 6--13 March) and, for every change of
the more active contract of a sibling pair, counting the changes of the quieter one at lags of
$-15$ to $+15$ minutes against the count expected under independence, two contracts of the same
game change in the same minute 2.7 times more often than chance, 2.3 times at one minute of
distance and still 1.9 times at fifteen; the contracts of the same election or ladder, 2.0 and
1.5 times.  But the histogram is symmetric to the third decimal (asymmetry between $+k$ and
$-k$ below 0.01 at every lag): the quieter contract does not follow the active one any more
than it precedes it.  The common cause, the event on the pitch or the headline, reaches both
books within the same minute, which is below the resolution of any horizon this thesis uses.
What a sibling did in the last hour is therefore already in the market's own book by the time
the model reads it, and the sibling features of \cref{sec:graph:siblings} lose rather than
gain in the trees.  The sign structure is weak and, for exclusive events, slightly
concordant rather than opposite (43--47\% of the simultaneous changes have opposite signs): a
part of the co-movement is a market maker re-quoting the same side of several books, not
probability mass moving from one candidate to another.

\begin{figure}[t]
\centering
\begin{tikzpicture}
\begin{axis}[thesis, xlabel={lag of the quieter contract's change, minutes (negative: it moved first)},
  ylabel={co-changes / expected under independence}, xmin=-15.5, xmax=15.5, ymin=0.9, ymax=3.3,
  legend pos=north east, xtick={-15,-10,-5,0,5,10,15}]
\addplot[vizA, mark=*, mark size=1.3pt] coordinates {(-15,1.99) (-14,2.0) (-13,2.01) (-12,2.03) (-11,2.04) (-10,2.06) (-9,2.07) (-8,2.08) (-7,2.11) (-6,2.12) (-5,2.15) (-4,2.18) (-3,2.21) (-2,2.25) (-1,2.32) (0,2.67) (1,2.33) (2,2.25) (3,2.2) (4,2.18) (5,2.15) (6,2.12) (7,2.1) (8,2.07) (9,2.05) (10,2.04) (11,2.02) (12,1.99) (13,1.98) (14,1.96) (15,1.94)};
\addplot[vizB, mark=*, mark size=1.3pt] coordinates {(-15,1.31) (-14,1.32) (-13,1.31) (-12,1.32) (-11,1.33) (-10,1.33) (-9,1.34) (-8,1.35) (-7,1.36) (-6,1.36) (-5,1.37) (-4,1.4) (-3,1.41) (-2,1.43) (-1,1.53) (0,1.99) (1,1.52) (2,1.43) (3,1.41) (4,1.4) (5,1.35) (6,1.35) (7,1.35) (8,1.34) (9,1.32) (10,1.32) (11,1.32) (12,1.31) (13,1.3) (14,1.3) (15,1.29)};
\addplot[vizGrid, dashed, line width=0.8pt] coordinates {(-15.5,1) (15.5,1)};
\legend{same game (160{,}662 pairs), same template and end date (18{,}099 pairs)}
\end{axis}
\end{tikzpicture}

\caption{Lead--lag between sibling contracts, 6--13 March: for every minute change of the
more active contract of a pair, the number of changes of the quieter one at each lag,
divided by the number expected under independence.  The peak is at lag zero and the two
sides are symmetric: siblings co-move, nobody leads.}
\label{fig:expl:leadlag}
\end{figure}

\subsubsection{Wallets}
\label{sec:expl:wallets}

Each day 122{,}000--161{,}000 distinct wallets trade.  A wallet's past accuracy, the share
of its trades on the side that the price then moved to over the horizon, is strongly
persistent: the Spearman correlation between a wallet's hit rate on the training days and
on the validation day is $0.5$--$0.6$.  Skill exists at the wallet level.  Whether it is
information about a \emph{market} is another question: a market is traded by many wallets,
their aggregate at any moment is close to that of random wallets with the same activity, and
as a market-level predictor the skill of the wallets that traded in the last hour adds about
one point over an identity-free placebo and does not beat Mean Reversion
(\cref{sec:graph:smartmoney}).  The trade of a skilled wallet is also a book event, seen by
the order-flow features before any wallet feature can see it.

